\clearpage
\part*{Part XIII. Thermodynamics}
\label{atlaspart:13-thermodynamics}
\addcontentsline{toc}{part}{Part XIII. Thermodynamics}
\section*{Introduction to Part XIII}\label{narr:part:XI}
\addcontentsline{toc}{section}{Introduction to Part XIII}
Thermodynamics enters the construction through a count at a boundary. Part XIII reads the second law as the count of records at the boundary of a web under acts, the thermal state as a fixed point of the twirl at a cap, the ensembles as readings at caps, the order parameter as a quantity read at the stop condition of a saturation, the phase transition as a change of sink class, and the remainder and the arrow as what is left when the forward saturation has run. The Part has more theorems than most, because the counting face of chapter 1-3 and the saturation of chapter 1-5 give it its objects directly.

The order of the seven chapters is the order of the textbook, the second law first and the arrow last, but the objects beneath the names are two counts and one direction. The first count is the boundary count, the number of records at the interface of a cut, which chapter 11-1 shows to be non-decreasing under the acts. The second count is the register count of the whole common support, written W throughout the Part, which chapter 11-4 reads as the order parameter and chapter 11-7 reads as the entropy. The direction is the forward saturation, the sequence of moves that joins pairs and never cuts them, along which both counts rise. The thermal state and the ensembles between them are what the twirl leaves fixed at a cap, and the phase transition is a change in the class of the sink a saturation reaches. A reader who holds those three objects in mind will find every box of the Part a statement about one of them, and will find the two open results of the Part, the boundary count from the valence law and the thermal state through a bath, at the two places where a second reading of the same object is admitted.

Chapter 11-1 reads the second law. Theorem~217.1 states the second law at the counting face and Theorem~217.2 the boundary count at the counting face, both theorems, with the boundary count from the valence law, Open result~217.3, open beside them. The object of the chapter is the boundary count of an interface at its counting face, which the literature calls the entanglement entropy across a cut, and the theorem states that this count is non-decreasing under every one-act read whose open end lies at the boundary and under every one-act read beside, at each of the thirteen bodies, Theorem~217.1. At every one of the sixty deciding cells of the horizon composite and the universe web the count at the counting face equals the site bound of the valence-law question at sixty of sixty and the free-direction bound at none, Theorem~217.2. The open result beside them finds that the boundary count the horizon results consume is not one bound of the counting face, equal to the site bound at three of four deciding elements, parting first at the closure of the universe web, and to the free-direction bound at two of four, Open result~217.3. Chapter 11-2 reads the thermal state as the twirl fixed point at a cap, Open result~218.1, and through the joined web with a bath body, Open result~218.2, both open. The thermal state is defined as the fixed point of the identification twirl at a cap at fixed boundary count, that fixed point is the flat state at one over the dimension of the cap, and the pair ket is re-met at every one of its twenty-four cells, Open result~218.1. Read instead through a joined network with a bath body, every bath of the direct-sum family joined to the one link at every pairable site over sixty-four joined networks, the marginal on the system caps equals the flat fixed point at the first order at every joined network and departs at the second order, which is the parting, Open result~218.2. Chapter 11-3 reads the ensembles as cap readings at the joined web in Theorem~219.1, a theorem that takes sixteen thermal names of the literature, the thermal state, the temperature, the zeroth and third laws, the partition function, the three ensembles, the three distributions, the equipartition, the chemical potential, the free energy, the thermodynamic limit and the heat capacity, and reads each as a reading at a cap of the joined web.

Chapter 11-4 reads the order parameter. Open result~220.1 states it at the stop condition as an open result, Theorem~220.2 reads it at the forward saturation and Theorem~220.3 is the distillation theorem at the forward saturation, both theorems. Over the configurations of the large closure, with its 44341 stable configurations, and of the clean twin, with its 7133 sinks, every web is classed by its open dock ends and the register count of the whole common support is read as an order parameter, Open result~220.1. Restricted to the forward moves, with the 2556 inverse edges excluded, that count is monotone along every forward edge, strictly increasing at the forward pairing edges and falling at the inverse edges alone, Theorem~220.2. At every forward edge of the twelve readable saturations the pair the move joins is, by the axiom, a maximally entangled pair of its common support, so its fidelity to the flat pair is one exactly at every link and monotone along every forward path, Theorem~220.3. Chapter 11-5 reads the phase transition as the sink-class change in Theorem~221.1, a theorem: the order parameter of the preceding chapter, the pair of the class of a sink and the register count, is re-met along every forward edge of the large closure at the per-web and order tables of the forward saturation, the count strictly increasing along the path and the class changing where the transition is read, Theorem~221.1. Chapter 11-6 holds the equilibrium theorem as awaited, the table titled The equilibrium theorem, awaited (in chapter~\ref{mod:XI-6-equilibrium-theorem}, among the instruments), and chapter 11-7 reads the thermodynamic remainder at the order parameter, Open result~223.1, open, and at the forward saturation, Theorem~223.2, a theorem, which is where the arrow of the construction stands. Over the move graph at the sinks of the large closure and the clean twin, the entropy production read as the change of the order parameter at every forward edge is non-decreasing at the 78785 forward edges and falling at the 1000 reverse edges, Open result~223.1, and at the forward saturation, with the register count taken as the entropy, the entropy production is strictly increasing at every forward edge of the large truncated closure, the second law holding at the strict reading along every forward path with the name entropy keyed, Theorem~223.2. The arrow of the construction is therefore the direction in which the forward saturation runs, and the reader should notice that it is read at the same count the Part began with.

The two counts are not the same object, and the Part keeps them apart deliberately. The boundary count is a property of an interface, a cut through a web, and it answers the question of how many records cross the cut; the register count is a property of the whole common support of a configuration, and it answers the question of how far the configuration has grown. The second law of chapter 11-1 is a statement about the first count at a fixed interface under reads, Theorem~217.1, while the arrow of chapter 11-7 is a statement about the second count along the forward edges of the saturation, Theorem~223.2. They meet at the horizon composite, where the boundary count at the counting face equals the site bound at every one of the sixty deciding cells, Theorem~217.2, and where the open result beside names the one place at which the two bounds of the counting face part, the closure of the universe web, Open result~217.3. For the relation between the two counts, the construction supplies the pair of counts at every cell and the direction in which each moves. The sixteen thermal names of chapter 11-3 are read in the same spirit, at the cap of the joined web and not at a limit, which is why the thermodynamic limit stands among them as a name read at a cap rather than as a limit taken, Theorem~219.1.

In the language of the literature, this Part names entropy, the second law, the thermal state, the Gibbs ensembles, the order parameter, the phase transition, equilibrium and the thermodynamic arrow. The correspondence is by name at each. The boxes establish counts at a boundary and fixed points of a twirl, and the arrow is read as the direction of the forward saturation.

The dependency lines of the Part point to the records, the moves, the saturation and the probe of Part I and to one another in order, the thermal state resting on the second law, the ensembles on both, and the transition on the thermal state and the order parameter. The second law itself rests on the counting face of chapter 1-3 and on the horizon composite of Part VI, which is where the sixty deciding cells are built, and the Gibbs state of chapter 8-5 is the same fixed point read at a cap of a body. A reader following the derivation graph will find Part XIII a compact chain hanging from the counting face, and will find that the one word the Part borrows most heavily, entropy, names a count that the construction had already been keeping since its first chapter.
\section{The second law as the boundary count under acts}\label{mod:XI-1-second-law-as-the-boundary-count-under-acts}
\chdesig{13-1}{11-1}{XI-1-second-law-as-the-boundary-count-under-acts}{XI-1-second-law-as-the-boundary-count-under-acts}
\label{narr:open:11-1}The second law is a count at a boundary. Theorem~217.1 states it at the counting face and Theorem~217.2 the boundary count there, both theorems, with the boundary count from the valence law, Open result~217.3, open beside them.
\begin{chtheorem}[The second law at the counting face]\label{box:XI-1-second-law-as-the-boundary-count-under-acts:1}
The object of this chapter is the boundary count (the entanglement entropy across a cut) of an interface at its counting face. The theorem states that this count is non-decreasing under every one-act read whose open end lies at the boundary, and under every one-act read beside, at each of the thirteen bodies and the neutral body: at all 228 one-act reads of the fourteen bodies (two per link at 114 links, 78 with an open far end, 150 washed) the count rises or stays at the site bound, at the free-direction bound and at the directions read, no decreasing act is found, and the running count along the order coordinate is non-decreasing. The result is a corollary of the record clause (a read is a further record and no read removes one) and of the definition of an open end, so the second law is a theorem at the counting face with the arrow of the order-coordinate theorem as its direction, the words entropy and arrow entering by name.
\end{chtheorem}
\chstanding{theorem}
\chdeps{Definitions 1.1 to 1.6, Theorems 18.1 to 18.4 and Open results 18.5 and 18.6 and Open results 241.1 and 241.2, Theorems 241.3 and 241.4 and Open results 241.5 to 241.7}{Open results 218.1 and 218.2, Theorem 219.1, Open result 220.1 and Theorems 220.2 and 220.3 and Open result 223.1 and Theorem 223.2}
A one-act read appends a record and removes none, and its far end is left unwashed exactly when the near end is the larger, so the count never falls along the order coordinate. The depth-monotonicity question is a corollary at the level of the records, where the depth decreases at none of the 32851 census moves, and stays open at the web, where the count of directed edges decreases at 22435 of them.
\begin{chtheorem}[The boundary count at the counting face]\label{box:XI-1-second-law-as-the-boundary-count-under-acts:2}
The parent theorem fixes the count: at every one of the 60 deciding cells of the horizon composite and the universe web the count at the counting face equals the site bound of the valence-law question at 60 of 60 and the free-direction bound at 0 of 60, so the first premise of the Einstein-shaped equation is a theorem of the count at the composite.
\end{chtheorem}
\chstanding{theorem}
\chdeps{Definitions 1.1 to 1.6, Theorems 18.1 to 18.4 and Open results 18.5 and 18.6 and Open results 241.1 and 241.2, Theorems 241.3 and 241.4 and Open results 241.5 to 241.7}{Open results 218.1 and 218.2, Theorem 219.1, Open result 220.1 and Theorems 220.2 and 220.3 and Open result 223.1 and Theorem 223.2}
The site bound, met at every deciding cell of the composite, is the bound whose rise the derivation below tabulates first.
\begin{chwall}[The boundary count from the valence law]\label{box:XI-1-second-law-as-the-boundary-count-under-acts:3}
Beneath both stands an open result: the boundary count the horizon results consume is not one bound of the counting face, equal to the site bound at 3 of 4 deciding elements (parting first at the closure of the universe web, 63 against 129) and to the free-direction bound at 2 of 4 (parting first at the horizon dock, 3 against 1).
\end{chwall}
\chstanding{open result}
\chdeps{Definitions 1.1 to 1.6, Theorems 18.1 to 18.4 and Open results 18.5 and 18.6 and Open results 241.1 and 241.2, Theorems 241.3 and 241.4 and Open results 241.5 to 241.7}{Open results 218.1 and 218.2, Theorem 219.1, Open result 220.1 and Theorems 220.2 and 220.3 and Open result 223.1 and Theorem 223.2}
At one face the count would be one bound at every deciding element. At the face this chapter inherits it is the site bound at three and the free-direction bound at two, so the second law is stated at both bounds and holds at each.
\begin{chbridge}
The literature calls the boundary count the entanglement entropy across a cut (nearest), the act an entanglement swap (nearest), a record an invariant fixed by projective measurements (nearest), and the result the second law, $\Delta S \geq 0$ for an isolated system \cite{Clausius1865}, with the H-theorem as its kinetic form \cite{Boltzmann1872}, both by name only.
\end{chbridge}
\begin{chproof}[census]
The premises are the Axiom and the six definitions of Section~\ref{mod:FRAMEWORK} (definitions), the record clause and the read transport of Section~\ref{mod:I-3-records-reads-and-the-born-count-the-gain-lemma} (a theorem), the move table of Section~\ref{mod:I-4-moves-in-one-table} (a table) and the parent theorem and the open result stated above. Nothing outside that list is used, and the chain from the axiom to this result is four chapters deep.
From the fixed library of networks (64 networks, 1334 sites, 1691 links), by the body list, we obtain the fourteen bodies carrying 114 links, and from the links, by the rule that a one-act read runs along a link in either orientation, the set counted, $N_{\mathrm{acts}} = 2 \cdot 114 = 228$. From the near and far dimensions of an act, by the read transport, we obtain the state of its far end, open when the near dimension exceeds the far and washed otherwise. Counting,
\begin{equation}
228 = 78 + 150, \qquad 78 = 24 + 54 ,
\end{equation}
the 78 open far ends falling 24 at a boundary site and 54 inside. From the sites of a body, by the definition of the boundary (a site at a degree below its dimension), we obtain the site bound as the sum of the dimensions of the boundary sites and the free-direction bound as the count of their free directions, at the electron body the boundary sites $a_3$, $w_1$ and $w_2$ of dimension two, so $B_{\mathrm{site}} = 2 + 2 + 2 = 6$ and $B_{\mathrm{free}} = 3$. From an act and the record clause (one pair in, one pair out), by the definition of an open end, we obtain the change of the bounds: at an open far end one unwashed direction is appended at the far site, the free bound rises by one and the site bound rises by the dimension of the far site when that site joins the boundary and by nothing when it stood there already, and at a washed end, as at the web face at every act, both stay,
\begin{equation}
\Delta B_{\mathrm{free}} = [\text{far end open}], \qquad \Delta B_{\mathrm{site}} = d_{\mathrm{far}} \, [\text{far end open and inside}] .
\end{equation}
From the 228 acts and the change rule, by tabulation, we obtain the counts by class, the free bound rising 78 (24 at the boundary, 54 inside), unmoved 150, falling 0, the site bound rising 54, unmoved 174, falling 0, the web face unmoved at 228 of 228, and the first decreasing act none. From the order coordinate (one step per arrow), by the running sum, we obtain the running count at 228 steps non-decreasing at the free face (rising 78) and at the site face (rising 103). The smallest instance is the electron body, 9 sites, 10 links, 20 acts, of which 14 are washed and move nothing, the acts from the hub $H$ to $A$, $a_1$ and $a_2$ (dimensions four to three, four to two, four to two) are open inside and raise the site bound by $3$, $2$ and $2$ and the free bound by one each, and the acts from $H$ to $a_3$, $B$ to $w_1$ and $C$ to $w_2$ are open at the boundary and raise the free bound alone, so
\begin{equation}
B_{\mathrm{site}} : 6 \to 13 \ (\text{rise } 3 + 2 + 2), \qquad B_{\mathrm{free}} : 3 \to 9 \ (\text{rise } 6),
\end{equation}
with no act lowering either. The property is met at every element, non-decreasing at the 24 boundary-open acts 24 of 24 at each bound and at all 228 acts 228 of 228 at each, which is the result.
A washed far end read as open moves the count (all 228 ends open under it). The record clause struck, so that an act removes a record, moves it, the first decreasing boundary-open act at the seventh link of the electron body. The valence face in place of the counting face moves the counts with the deciding cells unmoved. The order coordinate reversed reads non-increasing, the first decrease at the sixth link of the neutral body. The text struck moves every cell. The twelve controls hold.
The result was reproduced by an independent route, two engines written apart agreeing at every cell.
\chcert{228}
\end{chproof}

\par\medskip\noindent{\itshape The law read, the next chapter reads the state it leads to.}\label{narr:close:11-1}\par
\section{The thermal state}\label{mod:XI-2-thermal-state}
\chdesig{13-2}{11-2}{XI-2-thermal-state}{XI-2-thermal-state}
\label{narr:open:11-2}Open result~218.1 reads the thermal state as the twirl fixed point at a cap and Open result~218.2 the thermal state through the joined web with a bath body, both open results with their partings named.
\begin{chwall}[The thermal state as the twirl fixed point at a cap]\label{box:XI-2-thermal-state:1}
The thermal state is defined as the fixed point of the identification twirl (the wash) at a cap at fixed boundary count: that fixed point is the flat state at one over the dimension of the cap, and the maximally entangled pair state (the pair ket) is re-met at every one of its 24 cells. The definition does not reproduce the Gibbs form: at fixed count at the counting face the fixed point is not uniform over the records of that count at the one link, where at the count four one record stands at the weight one half and another at one third; at the fourteen bodies of the fixed library of networks it is uniform at every count at 2 of 14 bodies and the marginal is geometric at every cap at 1 of 14; the temperature remains a symbol with no value, and the alternative readings are printed with none chosen.
Re-posed at the washed-direction count with the budget replaced by the site bound six of the closure without caps: at the one (2,3) link the 172 records re-binned at the washed-direction count fall in the bins one, two, three and six at 79, 48, 42 and 3 records, the flat weight one over the count at every record, uniform at four of four bins by construction, while at the free-direction count beside the parting above re-meets at the count four, one record at one half against another at one third, typed at its cause, the two end dimensions of a read link; at the closure the fixed point is uniform at every bin at the washed count at 7698 of 7698 states, but the cap marginal at fixed washed count is geometric or vacuous at every site cap at 76 of 7698 states only, the first parting at the state of three qubits on a triangle at its first cap, the counts zero and one at the ratio 36 over 41 against one and two at 1 over 9, typed at its cause, the placements of the survived ends at the sites at one product; so the Gibbs form does not stand at the cap marginal at the washed count either, and the parting moves from the bin weights to the cap marginal, the temperature a symbol with no value as above.
\end{chwall}
\chstanding{open result}
\chdeps{Definitions 1.1 to 1.6, Theorems 28.1 and 28.2, Open results 28.3 and 28.4 and Theorem 28.5 and Theorems 217.1 and 217.2 and Open result 217.3}{Theorem 219.1, Open result 220.1 and Theorems 220.2 and 220.3 and Theorem 221.1}
At one face the Gibbs form would be an output of the fixed-point definition, the multiplier the base. At the face this chapter inherits the two caps of a read link carry two dimensions, so two records at one count carry two weights, and the parting is typed at its cause, the first at the count four where one record weighs one half and another one third.
\begin{chwall}[The thermal state through the joined web with a bath body]\label{box:XI-2-thermal-state:2}
Read instead through a joined network with a bath body (every bath of the direct-sum family joined to the one link at every pairable site, 64 joined networks), the marginal on the system caps equals the flat fixed point at the first order at every joined network and departs at the second order at every one of the 64, the first parting at the two-dock electron as the bath at its third site.
\end{chwall}
\chstanding{open result}
\chdeps{Definitions 1.1 to 1.6, Theorems 28.1 and 28.2, Open results 28.3 and 28.4 and Theorem 28.5 and Theorems 217.1 and 217.2 and Open result 217.3}{Theorem 219.1, Open result 220.1 and Theorems 220.2 and 220.3 and Theorem 221.1}
At one face the bath would leave the system at its flat fixed point at every order. At the face inherited here the marginal agrees at the first order and departs at the second at every joined network, the parting again the two end dimensions of a read link.
\begin{chbridge}
The literature calls the wash the twirl, the group average over the unfixed identifications (exact), the pair ket a maximally entangled pair (exact), the washed marginal a reduced density matrix with the twirl applied (exact), the cap a terminal site of a network (nearest), and the object sought the Gibbs state $\rho = e^{-\beta H}/Z$ of the canonical ensemble \cite{Gibbs1902} under the zeroth law \cite{Callen1985}, by name only.
\end{chbridge}
\begin{chproof}[census and exact cell]
The premises are the Axiom and the six definitions of Section~\ref{mod:FRAMEWORK} (definitions), the read at a register of Section~\ref{mod:I-6-read-at-a-register-the-probe-and-the-measurement} (a theorem) and the second law of Section~\ref{mod:XI-1-second-law-as-the-boundary-count-under-acts} (a theorem). Nothing outside that list is used, and the chain from the axiom to this result is five chapters deep.
\emph{The exact part.} The object is the one link of the (2,3) genre with its two caps and the pair space of its common support of dimension $d$. From the identification unitaries at a cap, by the wash, we obtain the twirl as a map on the $d^{4}$ matrix units of the pair space, the projection onto the commutant of the identification unitary paired with its conjugate, and at every support $d$ of the five (one, two, three, four and six) the map is idempotent, trace preserving and unital, fixes the pair projector and has exactly one pure fixed state, the pair ket, with a commutant and a fixed-point space of dimension one at the support one and of dimension two at the supports two, three, four and six. The smallest instance is the support two. The basis is the product basis of the two ends of the link, the four states $\lvert 00 \rangle$, $\lvert 01 \rangle$, $\lvert 10 \rangle$ and $\lvert 11 \rangle$ in this order, in which the matrix unit $E_{ij} \otimes E_{kl}$ carries $\lvert jl \rangle$ to $\lvert ik \rangle$ and the twirl of a unit is the orthogonal projection of that unit onto the span of the pair projector $P$ and its complement $\mathbf{1} - P$, so that the coefficient of $P$ is the trace of $P$ against the unit and the coefficient of $\mathbf{1} - P$ is the trace of $\mathbf{1} - P$ against the unit divided by $d^{2} - 1$. Writing $P$ for the pair projector onto the state
\begin{equation}
\lvert \Phi \rangle = \tfrac{\sqrt{2}}{2}\,(1, 0, 0, 1),
\end{equation}
the sixteen matrix units $E_{ij} \otimes E_{kl}$ are carried to
\begin{equation}
T(E_{ij} \otimes E_{kl}) = a \, P + b \, (\mathbf{1} - P) ,
\end{equation}
with $(a, b) = (0, 0)$ at ten units, $(0, \tfrac{1}{3})$ at two, $(\tfrac{1}{2}, -\tfrac{1}{6})$ at two and $(\tfrac{1}{2}, \tfrac{1}{6})$ at two. From the fixed-point space, by the flat weight, we obtain the weight one over the dimension at every cap,
\begin{equation}
w_{\mathrm{cap}} = \tfrac{1}{d}, \qquad d \in \{2, 3\} \text{ at the two caps of the one link},
\end{equation}
which is the thermal state of the definition.
\emph{The census part.} The set counted is the 172 records of the one link, generated by the closure at every order and typed to the sectors odd 79, even-mixed 48, even-same 42 and closure 3. From a record and the boundary count at the free face (one unwashed direction per survived end, the base counts the site bound five and the free-direction bound three), by the count law, we obtain its bin, and the bins hold 79 records at the count three, 90 at four and 3 at five. From a record, by the flat weight at every survived end, we obtain its fixed-point weight
\begin{equation}
w(R) = \prod_{\text{survived ends}} \tfrac{1}{d_{\mathrm{end}}}, \qquad w = \tfrac{1}{2} \text{ at a qubit end}, \quad w = \tfrac{1}{3} \text{ at a qutrit end},
\end{equation}
so the property tested, uniformity of the weight over the records of one count, holds at the bin three (the weight one at 79) and at the bin five (one sixth at 3) and fails at the bin four, where 48 records weigh one half and 42 weigh one third, uniform at 2 of 3 bins. The first parting element is the pair of records numbered 80 and 83 at the count four, at one half against one third. From the fixed point at fixed total count, by the marginal at a cap, we obtain at the cap $A$ the weights $0 \mapsto 1$ at the total three, $0 \mapsto \tfrac{7}{19}$ and $1 \mapsto \tfrac{12}{19}$ at the total four, and $1 \mapsto 1$ at the total five, geometric at every count of the one link because a cap of a single link admits the counts zero and one alone. At the fourteen bodies the same two tests give uniformity at every count at 2 of 14 bodies and a geometric marginal at every cap at 1 of 14. The first parting at the electron body is the count four, where $6\ 287\ 524\ 238\ 633\ 142\ 878\ 208$ records weigh one quarter against $13\ 420\ 516\ 941\ 450\ 723\ 557\ 376$ at one third, and its cap $A$ at the total five, where the counts zero and one stand at the ratio $36581654483132/129820737024193$ against the counts one and two at $450395748/10524066307$.
The check is by the counter-tests. The twirl struck, the raw state in place of the fixed point, moves the weights and reads every bin uniform. The boundary count replaced by the record count moves the bins to one bin of 172. The marginal taken at the whole web in place of the cap moves. A washed end counted as open moves the open ends from 78 to 228. The text struck moves every cell. The twelve controls hold, the second-law count non-decreasing at 228 of 228 acts and the pair ket re-met at 24 of 24 cells among them.
The result was reproduced by an independent route, two engines written apart agreeing at every cell. The second open result is stated in its own box.
\chcert{172}
\end{chproof}
\begin{figure}[h]
\begin{center}
\setlength{\unitlength}{1pt}
\begin{picture}(220,60)(0,0)
\put(40,30){\circle{24}}
\put(180,30){\circle{24}}
\put(52,30){\line(1,0){116}}
\put(35,26){$2$}
\put(175,26){$3$}
\put(8,4){cap $A$}
\put(170,4){cap $B$}
\put(85,36){the one link}
\end{picture}
\end{center}
\caption{The one link of the (2,3) genre. A qubit site and a qutrit site joined by one maximally entangled pair, with a cap at each end. The twirl acts at a cap on the pair space of the common support.}
\end{figure}

\par\medskip\noindent{\itshape The thermal state in hand, the next chapter reads the ensembles.}\label{narr:close:11-2}\par
\section{The ensembles as cap readings}\label{mod:XI-3-ensembles-as-cap-readings}
\chdesig{13-3}{11-3}{XI-3-ensembles-as-cap-readings}{XI-3-ensembles-as-cap-readings}
\label{narr:open:11-3}Theorem~219.1 reads the ensembles as cap readings at the joined web, a theorem resting on the second law and the thermal state.
\begin{chtheorem}[The ensembles as cap readings at the joined web]\label{box:XI-3-ensembles-as-cap-readings:1}
Sixteen thermal names of the literature (the thermal state, the temperature, the zeroth and third laws, the partition function, the three ensembles, the three distributions, the equipartition, the chemical potential, the free energy, the thermodynamic limit and the heat capacity) are read as readings taken at the caps of the joined networks with a bath body: the three ensembles, the partition sums, the bath readout and the zeroth and third laws are exact objects at the joined networks, the three distributions follow the sign of the own-site rotation with every parting named, and no Gibbs ratio is asserted. Thirteen names are derived and three are glossed, none open; a second reading at zero freedom, reading every derived name against the textbook object it stands for, typed five as proof sketches, where the textbook definition is instantiated at the established object, and eleven as namings, where an energy, a temperature, an exponential form, an occupation number, a particle number or an entropy at a least count is absent and the name attaches by resemblance, both typings printed and the deciding letter the same at both.
\end{chtheorem}
\chstanding{theorem}
\chdeps{Definitions 1.1 to 1.6, Theorems 217.1 and 217.2 and Open result 217.3 and Open results 218.1 and 218.2}{no theorem beyond it at this release}
The three glosses each name the clause the construction lacks. An energy object at the free energy, where a count of the register exponent stands in place of a function of an energy and a temperature. The limit at an unbounded copy count at the thermodynamic limit, every readout exact at two, three and four copies and defined at no larger count, and a temperature object at the heat capacity, the temperature a symbol with no value. The two open results of the preceding chapter are the objects the readings are taken at.
\begin{chbridge}
The literature calls the washed marginal a reduced density matrix with the twirl applied (exact), the joined web the composition of two tensor networks along shared sites (nearest), the cap a terminal site (nearest), and the objects named the ensembles of Gibbs with the partition function $Z = \sum_n e^{-\beta E_n}$, the free energy $F = -\beta^{-1}\ln Z$ and the grand sum $\Xi$ \cite{Gibbs1902}, with the distributions, the equipartition and the heat capacity $C = \partial U/\partial T$ \cite{Callen1985}, by name only.
\end{chbridge}
\begin{chproof}[census and exact cell]
The premises are the Axiom and the six definitions of Section~\ref{mod:FRAMEWORK} (definitions), the second law of Section~\ref{mod:XI-1-second-law-as-the-boundary-count-under-acts} (a theorem) and the two open results of the thermal state at Section~\ref{mod:XI-2-thermal-state} (open results, read at their faces). Nothing outside that list is used, and the chain from the axiom to this result is six chapters deep.
\emph{The exact part.} The object is the one (2,3) link joined at every pairable site to a bath body of the direct-sum family, 64 joined networks, every record weighing its act readout. The marginal at the two system caps $X$ and $Y$ is read order by order. From the flat fixed point of the preceding chapter, by the marginal at the order zero, we obtain the microcanonical reading, the bin of the one link at the count three, 79 records at the weight one, uniform, with the cap marginals
\begin{equation}
A: \ 0 \mapsto \tfrac{7}{19}, \ 1 \mapsto \tfrac{12}{19}; \qquad B: \ 0 \mapsto \tfrac{12}{19}, \ 1 \mapsto \tfrac{7}{19},
\end{equation}
the order-zero records of every joined network equal to that bin at 64 of 64. From the bath records summed at every order, by the share weight, we obtain the canonical reading, 384 cells at 64 joined networks, every cell an exact element of the field, equal to the flat fixed point at the order zero at 64 of 64 and at the order one at 64 of 64, where every joined network reads
\begin{equation}
P^{(1)}_{X} = \tfrac{21}{37}, \qquad P^{(1)}_{Y} = \tfrac{16}{37},
\end{equation}
and parting at the order two at 0 of 64, the first parting at the two-dock electron as the bath at its third site, where the cap $X$ reads $\tfrac{6307}{11043}$ against the flat value $\tfrac{25851}{45131}$. From the bath joined at every one of its pairable sites at once, by the same sum, we obtain the grand canonical reading at 19 baths with a joined network, and from the normalizer of each marginal the partition sum, its logarithm the free energy by name. The smallest instance is the bath of the docked up quark with the pairable set of two sites, the partition sums $Z_{1} = 2940$ and the grand reading equal to the canonical at the order one and not at the order two.
\emph{The census part.} The set counted is the sixteen thermal names listed above. From each name, by the typing rule (derived if its kind is exact and every cell defined, glossed if its kind is a name and every cell defined, open otherwise with its gate named), we obtain the counts by class, derived 13, glossed 3, open 0, the first open none. The equality of the system marginals at the order one across two baths is total at 561 of 561 pairs of the library baths and 2016 of 2016 pairs of all baths and transitive, which is the zeroth law, and it parts at the order two at 23 of 561 and 137 of 2016. From the own-site sign of the baths, by the sign class, we obtain the three distributions. The sign-odd class (the Fermi form by name) at 8 baths and 22 joined networks, the sign-even class (the Bose form by name) at 1 bath and 1 joined network, and the untyped class at 10 baths and 41, the carriers at the pairable set none at 4 of 4. From the second reading at zero freedom, by the instantiation test, we obtain five proof sketches (the zeroth law, the thermal state, the canonical and microcanonical ensembles, the partition function) and eleven namings, the deciding count the same at both typings.
The check is by the counter-tests. The bath struck moves 43 of 200 compared cells and the canonical reading falls to the microcanonical. The own-site sign struck moves 27 of 66 and the Bose and Fermi classes merge. The pairable set closed moves 97 of 175 and the grand canonical falls to the canonical. The text struck leaves the gate unmet and builds no cell. The eight controls hold, the marginal table of the bath open result re-derived at 960 of 960 rows and the flat fixed point at 5 of 5 supports among them.
The result was reproduced by an independent route, two engines written apart agreeing at every cell of the 64 joined networks.
\chcert{16}
\end{chproof}
\begin{figure}[h]
\begin{center}
\setlength{\unitlength}{1pt}
\begin{picture}(260,70)(0,0)
\put(40,35){\circle{22}}
\put(120,35){\circle{22}}
\put(51,35){\line(1,0){58}}
\put(36,31){$X$}
\put(116,31){$Y$}
\put(200,59){\circle*{6}}
\put(200,11){\circle*{6}}
\put(230,35){\circle*{6}}
\put(200,59){\line(5,-4){30}}
\put(200,11){\line(5,4){30}}
\put(131,35){\line(5,3){40}}
\put(131,35){\line(5,-3){40}}
\put(171,59){\line(1,0){29}}
\put(171,11){\line(1,0){29}}
\put(60,52){the system link}
\put(190,0){the bath body}
\end{picture}
\end{center}
\caption{The joined web. The one link of the system at its two caps $X$ and $Y$, joined at its pairable sites to a body of the fixed library as the bath. The ensembles are readings of the marginal at the two system caps order by order.}
\end{figure}

\par\medskip\noindent{\itshape The next chapter reads the quantity that tells one phase from another.}\label{narr:close:11-3}\par
\section{The order parameter at the stop condition}\label{mod:XI-4-order-parameter-at-the-stop-condition}
\chdesig{13-4}{11-4}{XI-4-order-parameter-at-the-stop-condition}{XI-4-order-parameter-at-the-stop-condition}
\label{narr:open:11-4}Open result~220.1 states the order parameter at the stop condition as an open result, Theorem~220.2 reads it at the forward saturation and Theorem~220.3 is the distillation theorem at the forward saturation, both theorems.
\begin{chwall}[The order parameter at the stop condition]\label{box:XI-4-order-parameter-at-the-stop-condition:1}
Over the configurations of the large closure (101973 webs and 44341 stable configurations, the sinks, the closure truncated at its budget) and of the clean twin (16010 webs, 7133 sinks), every web is classed by its open dock ends, and the register count W of the whole common support, read as an order parameter, is not monotone along every root-to-sink path, the first non-monotone edge at a reverse move; so at the unrestricted stop condition the order parameter is an open result.
Re-posed at the closure without caps at the site bound six, 7698 states and 93654 steps certified in Lean: the register count of the whole common support re-derived at every state and read along every path from a seed to a forward sink at the full edges is not monotone, 46499 non-monotone edges among 93610, the first at wave zero by an inverse move from a state at four to a state at two; the four transition kinds the previous reading found absent are present at four of four at the pairing-move edges and at every edge; at the forward edges alone the count is not monotone either, the first falling edge at wave one by the hub-forming move from thirty-two to sixteen, and at the tree edges at wave two by an inverse move, the first parting above re-met and re-read, so the parting stands at the site bound at every reading with the absent kinds present.
\end{chwall}
\chstanding{open result}
\chdeps{Definitions 1.1 to 1.6, Theorems 217.1 and 217.2 and Open result 217.3, Open results 218.1 and 218.2 and Open results 245.1 and 245.2, Theorem 245.3, Open results 245.4 to 245.8 and Theorem 245.9}{Theorem 221.1 and Open result 223.1 and Theorem 223.2}
At one face the order parameter would be an output of the sink class with a transition of each of the four kinds named at a wave and the register count monotone to the sink. At the face this chapter inherits one kind of four is present and the count is not monotone, the parting named with both objects at the derivation below. The two theorems that follow restrict the paths to the forward moves and supersede the open result there.
\begin{chtheorem}[The order parameter at the forward saturation]\label{box:XI-4-order-parameter-at-the-stop-condition:2}
Restricted to the forward moves (182940 forward edges of 185496 at the large truncated closure, the 2556 inverse edges excluded), W is monotone along every forward edge, strictly increasing at the forward pairing edges and falling at the inverse edges alone; the class changes along the forward edges at 9377 edges at 5 of 12 kinds, the first from the charged class to the dark class at the first wave by the attach move at the genre (2,3), and no web changes its determined class along a forward path; so the order parameter is the pair of the class and W at every sink, a theorem at the forward saturation.
\end{chtheorem}
\chstanding{theorem}
\chdeps{Definitions 1.1 to 1.6, Theorems 217.1 and 217.2 and Open result 217.3, Open results 218.1 and 218.2 and Open results 245.1 and 245.2, Theorem 245.3, Open results 245.4 to 245.8 and Theorem 245.9}{Theorem 221.1 and Open result 223.1 and Theorem 223.2}
\begin{chtheorem}[The distillation theorem at the forward saturation]\label{box:XI-4-order-parameter-at-the-stop-condition:3}
At every forward edge of the twelve readable saturations (322516 webs, 138129 sinks, 570349 forward edges) the pair the move joins is, by the Axiom, a maximally entangled pair of its common support, so its fidelity to the flat pair is one exactly at every link, monotone along every forward path, every sink at full fidelity; the distillation protocols of the literature, each a rise of the fidelity from below full, are located as no path of the saturation.
\end{chtheorem}
\chstanding{theorem}
\chdeps{Definitions 1.1 to 1.6, Theorems 217.1 and 217.2 and Open result 217.3, Open results 218.1 and 218.2 and Open results 245.1 and 245.2, Theorem 245.3, Open results 245.4 to 245.8 and Theorem 245.9}{Theorem 221.1 and Open result 223.1 and Theorem 223.2}
Both theorems are compositions riding the lowest grade in their chains. The open results among their parents (the large closure unstopped at its budget, the move table, the move action, the two seeded saturations and the library seeds) are read as they stand.
\begin{chbridge}
The literature calls the register the tensor factors of the common support (nearest), the sink a stable configuration of the saturation (nearest), the class of a web its charge sector by the open ends (nearest), and the objects named an order parameter in the picture of Landau \cite{Landau1937} and the fidelity monotone of entanglement distillation \cite{Bennett1996}, by name only.
\end{chbridge}
\begin{chproof}[census]
The premises are the Axiom and the six definitions of Section~\ref{mod:FRAMEWORK} (definitions), the saturation of Section~\ref{mod:I-5-saturation-in-one-figure-seeds-waves-termination-sinks} (a theorem), the second law of Section~\ref{mod:XI-1-second-law-as-the-boundary-count-under-acts} (a theorem) and the thermal state of Section~\ref{mod:XI-2-thermal-state} (open results, read at their faces). Nothing outside that list is used, and the chain from the axiom to this result is six chapters deep.
The set counted is the large closure at the site maximum of the fixed library, read to its fourth wave and stopped at its budget there (so a closure never, a truncated closure), 101973 webs with 44341 sinks, and beside it the clean twin, 16010 webs with 7133 sinks. From a web, by the readability order (the open dock ends), we obtain its class, one of dark, charged, neutral and beyond-two, and the sink census
\begin{equation}
\text{dark } 38700, \quad \text{charged } 5455, \quad \text{neutral } 185, \quad \text{beyond-two } 1, \qquad 44341 \text{ sinks}.
\end{equation}
From the pairing-move edges, 167113 at the large truncated closure, by the class at the parent and at the child, we obtain the transitions of class, 9013 changes, dark to charged at 8545 and charged to neutral at 468 and every other kind at zero, so the four kinds of the item are present at 1 of 4 (the three absent being neutral to dark, dark to neutral and neutral to charged). The first charged-to-neutral transition runs from the web numbered 84 to the web numbered 2496 at the third wave. From a web, by the product of the dimensions of its links at the whole common support, we obtain the register count $W$ with its exponent pair, and along the 81209 edges of the root-to-sink paths
\begin{equation} \adjustbox{max width=\dimexpr\linewidth-3em\relax}{$\displaystyle \begin{gathered} W \text{ is not monotone at } 1062 \text{ edges}, \\ \text{the first: } W = 64 \ (\text{web } 5) \to W = 4 \ (\text{web } 16) \text{ at the second wave by a reverse move}, \end{gathered}$} \end{equation}
every other reading of the register (the sum, the cap and the flat cap) not monotone beside. The smallest instance is the three roots and the first wave. The root webs numbered 1 (charged, $W = 2$), 2 (dark, $W = 2$) and 3 (dark, $W = 8$), the first child the web numbered 4 (dark, $W = 4$) by the attach move at the genre (2,3) from the charged root, a change of class at $W$ from two to four. The order parameter at a sink is the pair
\begin{equation}
(\text{class}, W): \quad (\text{charged}, 2), \ (\text{dark}, 8), \ (\text{dark}, 4), \ (\text{dark}, 64), \ \ldots
\end{equation}
printed at every one of the 44341 sinks, the register at the sinks running from 1 to $2^{65}$ at 216 distinct values. The first parting element of the open result is therefore the pair of the absent kinds and the first non-monotone edge named above. The forward restriction of the first theorem removes the 2556 inverse edges and leaves $W$ strictly increasing at the 182940 forward edges, read here at its statement.
The check is by the counter-tests. The readability order reversed (charged at two open dock ends, neutral at one) moves the classes and the deciding cells. The register count read at the body alone, every link carrying an open end struck from the product, moves the register at the sinks. The text struck leaves the gate unmet and every cell struck. Sixteen controls were posed and fourteen hold, the two that differ a control of the count of wave rows against the waves with webs, a count of two different sets and no identity, so the control cell reads as printed and the result stands at its face.
The result was reproduced by an independent route, two engines written apart agreeing at every class and every register of the 101973 webs.
\chcert{101973}
\end{chproof}

\par\medskip\noindent{\itshape The order parameter read, the next chapter reads the transition it marks.}\label{narr:close:11-4}\par
\section{The phase transition as the sink-class change}\label{mod:XI-5-phase-transition-as-the-sink-class-change}
\chdesig{13-5}{11-5}{XI-5-phase-transition-as-the-sink-class-change}{XI-5-phase-transition-as-the-sink-class-change}
\label{narr:open:11-5}Theorem~221.1 reads the phase transition as the sink-class change, a theorem resting on the saturation, the thermal state and the order parameter.
\begin{chtheorem}[The phase transition as the sink-class change]\label{box:XI-5-phase-transition-as-the-sink-class-change:1}
The order parameter of the preceding chapter, the pair of the class of a sink and the register count W, is re-met along every forward edge of the large closure (101973 webs, 44341 sinks, 182940 forward edges) at the per-web and order tables of the forward saturation, W strictly increasing along the forward edges, and tabulated beside at the proton seeded saturation (218340 webs), the bound-respecting seeded saturation (93481 webs, W unread at its table) and the library seeds (64 seeds at 1203 sink entries). The first critical W is found: at the large closure the class changes at 9377 forward edges (9270 at the fourth wave, 107 below it), the critical edge running from a charged parent at W two to a dark child at W four at the first wave by the attach move at the genre (2,3), so the critical W is four with the charged and the dark classes on either side; at the proton seed the critical W is 93312 (from the class beyond two to the neutral class), at the library seeds 5184 (charged to dark), and at the bound-respecting seed the critical measure is 21, each closure its own. The thermal state is located: fourteen bodies rebuilt from the caps stand eleven at a class boundary and none below the critical W of the large truncated closure, and the ratio two of the critical edge equals the twirl at the support at the dimension two alone. The identification with the phase transition of the literature is a correspondence by name, and every count is a count of a truncated closure.
\end{chtheorem}
\chstanding{theorem}
\chdeps{Definitions 1.1 to 1.6, Open results 218.1 and 218.2, Open result 220.1 and Theorems 220.2 and 220.3 and Open results 245.1 and 245.2, Theorem 245.3, Open results 245.4 to 245.8 and Theorem 245.9}{no theorem beyond it at this release}
The theorem is a composition riding the lowest grade in its chain, conditional on nothing beyond the text and the proved results it rests on. The open results among its parents, the order parameter at the stop condition and the thermal state among them, are read as they stand.
\begin{chbridge}
The literature calls the sink a stable configuration of the saturation (nearest), the class of a web its charge sector by the open ends (nearest), the register the tensor factors of the common support (nearest), and the object named a phase transition in the picture of Landau with its order parameter $\eta$ \cite{Landau1937}, the critical point of the two-dimensional Ising model its exactly solved case \cite{Onsager1944}, by name only.
\end{chbridge}
\begin{chproof}[census]
The premises are the Axiom and the six definitions of Section~\ref{mod:FRAMEWORK} (definitions), the saturation of Section~\ref{mod:I-5-saturation-in-one-figure-seeds-waves-termination-sinks} (a theorem), the thermal state of Section~\ref{mod:XI-2-thermal-state} (open results, read at their faces) and the order parameter of Section~\ref{mod:XI-4-order-parameter-at-the-stop-condition} (a theorem). Nothing outside that list is used, and the chain from the axiom to this result is seven chapters deep.
The set counted is the large truncated closure (101973 webs, 44341 sinks, 185496 edges of which 182940 are forward and 2556 inverse), with the proton seeded saturation (218340 webs), the bound-respecting seeded saturation (93481 webs) and the 64 library seeds beside. From a web, by the readability order and the product of the link dimensions, we obtain its class and its register count $W$ with its exponent pair, equal to the per-web and order tables of the forward saturation at every one of the 101973 webs and 44341 sinks, $W$ strictly increasing along the forward edges with no falling and no equal pair. From a forward edge, by the class of its parent and child, we obtain a transition where the class of the child differs, 9377 transitions at the large truncated closure, 9270 at the fourth wave and 107 below it. From the first change along every forward path, by its register, we obtain the critical edge,
\begin{equation}
\text{web } 1 \ (\text{charged}, W = 2) \ \xrightarrow{\ \text{attach at } (2,3)\ } \ \text{web } 4 \ (\text{dark}, W = 4), \qquad W_{\mathrm{child}} / W_{\mathrm{parent}} = 2 ,
\end{equation}
at the first wave, so $W_{c} = 4$, the exponent pair of $W_{c}$ being $(2, 0)$ in the dimensions two and three, the distinct registers at the first changes numbering 142, the least 4, 8, 12, 24 and 36. From the same rule at the other closures, by their own tables, we obtain each closure its own critical count, at the proton seed 15328 transitions and the critical edge from the root (beyond two, $W = 46656$) to the web numbered 5 (neutral, $W = 93312$) at the first wave by the same attach move, at the library seeds the critical $W = 5184$ from charged to dark, and at the bound-respecting seed the critical measure 21, its $W$ unread at its table. From the fourteen bodies rebuilt from the caps, by their class and register, we obtain their location against the critical $W$ of every closure, at the large truncated closure eleven at a class boundary and none below $W_{c}$, the smallest instance the docked up quark at $W = 4$ exactly, the neutral class, at a class boundary, and the electron body at $W = 5184 = 2^{6} \cdot 3^{4}$, above $W_{c}$ at the large truncated closure and at a class boundary at the library seeds, while the twirl flat weights at the five supports (1, 1/2, 1/3, 1/4, 1/6) meet the ratio two of the critical edge at the support two alone. The property tested, a change of the made class at a definite $W$ along the forward edges, holds at the large truncated closure and not only at its cap, which is the result, and the name of the phase transition is keyable at the critical $W$.
The check is by the counter-tests. The forward and inverse partition reversed moves the transitions, the critical $W$ and the deciding cells. The register read at the body alone, every link carrying an open end struck from the product, moves the critical $W$ at every closure. The text struck leaves the gate unmet. Eighteen controls hold, the per-web and order tables re-met at every one of the 101973 webs and the 44341 sinks and the bodies table at 14 of 14 among them, one sub-clause of one control re-spelled at review as a count of two sets and no identity.
The result was reproduced by an independent route, two engines written apart agreeing at every cell.
\chcert{101973}
\end{chproof}

\par\medskip\noindent{\itshape The equilibrium theorem is awaited in the next chapter.}\label{narr:close:11-5}\par
\section{The instrument tables of this Part}\label{mod:XIII-L-the-instrument-tables-at-the-repository}
\chdesig{13-L}{}{XIII-L-the-instrument-tables-at-the-repository}{XIII-L-the-instrument-tables-at-the-repository}
\begin{chtable}[The instrument tables of this Part]\label{box:XIII-L-the-instrument-tables-at-the-repository:1}
This Part lists 1 instrument chapter, printed among the instruments at the end of this compendium as chapter~\ref{mod:XI-6-equilibrium-theorem}, and the table beneath names each one with its key cells.
\end{chtable}
\chstanding{instrument}
\chdeps{Definitions 1.1 to 1.6}{no theorem beyond it at this release}
The listing moves no cell and states no result of its own. Each instrument stands at the repository as its chapter stood at the previous release of this compendium, and its tables are read there.
\begin{chbridge}
The tables listed here are instruments of the compendium, and this listing draws no correspondence with the literature.
\end{chbridge}
\begin{chproof}[instrument]
The listing in full. For each instrument chapter, printed among the instruments at the end of this compendium, the table gives its title, the result boxes it carries, the standing of each box with its citation, and the count of objects its certificate sentence names where it prints one.
{\small
\begin{longtable}{p{38mm}p{63mm}p{30mm}p{14mm}}
\hline
chapter & results & standing & objects \\
\hline
\endhead
The equilibrium theorem & The equilibrium theorem, awaited & instrument & -- \\
\hline
\end{longtable}
}
The rule that produced it. Every chapter of this Part whose first standing beyond a placeholder reads instrument is printed among the instruments at the end of this compendium with its tables, and its key cells are the titles of its boxes, its standing lines and its certificate count, read from its text as it stood at the previous release.
No derivation stands in this chapter; the tables stand at the repository.
\emph{Check.} Every file listed was re-read at the repository and read the same as at the previous release, and the listing was re-cut from the files by a second route.
Nothing stands open at this listing. At this release three results of this Part are re-graded at their chapters: the thermal state (the chapter 13-2) re-posed at the washed-direction count, the order parameter (the chapter 13-4) re-posed at the site bound, the thermodynamic remainder (the chapter 13-7) re-posed at the first missing clause.
\end{chproof}
\label{chend:XIII-L-the-instrument-tables-at-the-repository}
\section{The thermodynamic remainder and the arrow}\label{mod:XI-7-thermodynamic-remainder-and-the-arrow}
\chdesig{13-7}{11-7}{XI-7-thermodynamic-remainder-and-the-arrow}{XI-7-thermodynamic-remainder-and-the-arrow}
\label{narr:open:11-7}Open result~223.1 reads the thermodynamic remainder at the order parameter as an open result and Theorem~223.2 at the forward saturation as a theorem, where the arrow of the construction stands.
\begin{chwall}[The thermodynamic remainder at the order parameter]\label{box:XI-7-thermodynamic-remainder-and-the-arrow:1}
Over the move graph at the sinks of the large closure (a truncated closure) and the clean twin, the entropy production read as the change of the order parameter W at every forward edge is non-decreasing at the 78785 forward edges and falling at the 1000 reverse edges, while the boundary count of the second law parts from it at the first wave; the fluctuation relation read as the ratio of the forward to the reverse count at every pair of a move kind and its inverse is defined at 16 of 111 pairs; and the Onsager symmetry read as the one-act readout of the wash at the two ends of every link of the joined networks is equal at 19 of 43 at the full reading and at every pair at the common support; so at the unrestricted stop condition the entropy production is the first name with a missing clause, an open result.
Re-posed over the certificate graph of the closure without caps at the site bound six, 7698 states, 120149 edges and forty-one forward sinks, with the acts oriented by their direction: the entropy production read as the change of the register count at every forward edge is not non-decreasing at the 64951 forward edges, the first falling edge by the hub-forming act from sixteen to eight, strictly increasing at the 20705 pairing edges, and the boundary count of the second law compared at the same edges parts at its first edge; the fluctuation relation at the 172 pairs of an act and its inverse is given by no one function of the production, the first failing pair kind at the ratio 1742 over 1595 against 234 over 209 at the same production; and the Onsager readouts at the two joined webs, neither within the bound, are asymmetric at the ends and at the full support and equal at the common support; so the first name with a missing clause at the site bound is the entropy production, the register count falling at a forward edge, a clause other than the one the previous reading named, and the parting stands with its clause moved.
\end{chwall}
\chstanding{open result}
\chdeps{Definitions 1.1 to 1.6, Theorems 217.1 and 217.2 and Open result 217.3 and Open result 220.1 and Theorems 220.2 and 220.3}{no theorem beyond it at this release}
At one face the three names, the entropy production, the fluctuation relation and the Onsager symmetry, would be keyed with no clause missing. At the face this chapter inherits the first name parts at the second-law boundary count against the register count at a forward edge, the parting named with both objects at the derivation below.
\begin{chtheorem}[The thermodynamic remainder at the forward saturation]\label{box:XI-7-thermodynamic-remainder-and-the-arrow:2}
At the forward saturation, with W taken as the entropy, the entropy production is strictly increasing at every one of the 182940 forward edges of the large truncated closure, the second law holds at the strict reading along every forward path and the name entropy is keyed; the boundary count of the second law, read beside at the same edges, is typed two objects and not the entropy; the fluctuation relation has 179 defined pairs of 2989 at the sinks graph, the two candidate functions meeting the ratio at 0 and at 6 of them, so the fluctuation name stays silent, a count of a truncated closure with no ensemble; and the asymmetric Onsager pairs are exactly the pairs at a non-zero production difference at every one of the 43, so the Onsager name is keyed.
\end{chtheorem}
\chstanding{theorem}
\chdeps{Definitions 1.1 to 1.6, Theorems 217.1 and 217.2 and Open result 217.3 and Open result 220.1 and Theorems 220.2 and 220.3}{no theorem beyond it at this release}
The theorem is a composition riding the lowest grade in its chain. The open results among its parents are read as they stand, and the arrow enters by name with the second law of the first chapter of this Part as its direction.
\begin{chbridge}
The literature calls the register the tensor factors of the common support (nearest), the move an elementary operation on the network of pairs (nearest), the wash the twirl (exact), the one-act readout the outcome probabilities of a projective measurement (exact), and the objects named the entropy production of an irreversible process, the fluctuation relations of Crooks and Jarzynski \cite{Jarzynski1997,Crooks1999} and the Onsager reciprocity $L_{ij} = L_{ji}$ \cite{Onsager1931}, by name only.
\end{chbridge}
\begin{chproof}[census]
The premises are the Axiom and the six definitions of Section~\ref{mod:FRAMEWORK} (definitions), the second law of Section~\ref{mod:XI-1-second-law-as-the-boundary-count-under-acts} (a theorem) and the order parameter of Section~\ref{mod:XI-4-order-parameter-at-the-stop-condition} (an open result and two theorems, read at their faces). Nothing outside that list is used, and the chain from the axiom to this result is seven chapters deep.
The set counted is the move graph at the sinks of the large truncated closure (101973 webs, 44341 sinks), 79785 edges of which 78785 are forward and 1000 reverse, with the clean twin beside. From an edge, by the register count $W$ at its parent and its child as an exponent pair in the dimensions two and three, we obtain the production of the edge, read at the exponent pair, with the edge from web 1 to web 4 at the first wave as the example,
\begin{equation}
s(e) = W(\text{child}) - W(\text{parent}), \quad (1,0) \to (2,0), \ s = (1,0), \ W_{\mathrm{child}} / W_{\mathrm{parent}} = 2 ,
\end{equation}
and at the forward edges $W$ is non-decreasing (zero falling, the first none), strictly increasing at the 70942 pairing edges and falling at every one of the 1000 reverse edges, the entropy production read at $W$. From the second-law boundary count at the site face at the same edges, by comparing the two falls, we obtain the parting of the first name: at 3840 forward edges the boundary count falls while $W$ rises or rises while $W$ falls, the first at the edge from the root web 1 (site bound 5, $W = 2$) to the web numbered 8 (site bound 3, $W = 96$) at the first wave by the hub-step move with the kept leg, so the name entropy production is unkeyed, the second-law count parting from the register count. From every pair of a forward move kind and its inverse, by the forward and the reverse edge counts at a production, we obtain the fluctuation ratio at 111 pairs, defined at 16 and undefined at 95 at a zero reverse count, the first failing pair the attach at the genre (2,3) at the production $(1,0)$ with 17 forward and 0 reverse edges. At the 16 defined pairs the two candidates, the ratio of the registers and the unit, meet the ratio at 0 of 16 each, the smallest defined instance the cut move at
\begin{equation}
N_{\mathrm{forward}} / N_{\mathrm{reverse}} = 2136 / 536 = 267 / 67 .
\end{equation}
From the one-act readout at the two ends of every link of the two joined networks at the copy count two (the pair and the double pair, 43 end pairs), by comparing the two directions as exact objects, we obtain the Onsager reading, equal at 19 of 43 at the full readout, asymmetric, the first asymmetric pair the fourth link of the pair network at the genre (3,4), the hub-to-$A$ readout $(1/3, 1/3, 1/3)$ against the $A$-to-hub readout $(1/3, 1/3, 1/3, 0)$, while on the common support the two agree at 43 of 43, so the Onsager name is unkeyed at the full reading and equal at the support. The first parting element of the open result is therefore the entropy production at the edge named above, and the forward restriction of the theorem keys the entropy and the Onsager names and leaves the fluctuation name silent, read here at its statement.
The check is by the counter-tests. The orientation of the move graph reversed moves the productions, the letters and the first missing clause. The register read at the body alone, every link between sites of unequal dimension struck from the exponent pair, moves the second-law parting and the first failing pair. The text struck leaves the gate unmet and every cell struck. Seventeen controls were posed and sixteen hold, the one that differs a count of two sets and no identity, so the control cell reads as printed and the result stands at its face.
The result was reproduced by an independent route, two engines written apart agreeing at every cell.
\chcert{79785}
\end{chproof}

\par\medskip\noindent{\itshape Heat and the arrow read, Part XIV turns to lattices and atoms.}\label{narr:close:11-7}\par
